\documentclass[11pt]{article}
\input{style-sujet}

\newcommand{\note}[1]{\par\smallskip{\itshape\small #1}\par\smallskip}

\begin{document}

\entetesujet{Sujet bac 2019 -- Série C}

% ====================== EXERCICE 1 ======================
\exo{1}{4 points}

On considère l'équation $(E) : 109x - 226y = 1$ où $x$ et $y$ sont des entiers relatifs.

\begin{enumerate}
  \item Déterminer le PGCD de 109 et 226.
  \item Que peut-on conclure pour l'équation $(E)$ ?
  \item Vérifier que 141 est l'inverse de 109 modulo 226.
  \item \begin{enumerate}[label=\alph*.]
    \item Montrer que l'équation $(E)$ est équivalente à l'équation $(E') : 109x \equiv 1\ [226]$.
    \item Résoudre l'équation $(E')$.
    \item En déduire la solution de $(E)$.
    \note{Modifier la question : « En déduire les solutions de $(E)$. »}
  \end{enumerate}
\end{enumerate}

% ====================== EXERCICE 2 ======================
\exo{2}{8 points}

Dans le plan orienté, on considère un carré $ABCE$ de sens direct, de centre $O$.

$D$ est le symétrique orthogonal du point $O$ par rapport à la droite $(BC)$. $K$ et $O'$ sont les milieux respectifs des segments $[BD]$ et $[BC]$. On désigne par $G$ le symétrique du point $C$ par rapport à $B$.

\begin{enumerate}
  \item Faire une figure. On disposera $[BC]$ horizontalement, et on prendra $BC = 6$ cm.
  \item Soit $f$ la rotation qui transforme $O$ en $D$ et $A$ en $C$.
  \begin{enumerate}[label=\alph*.]
    \item Déterminer le centre de la rotation $f$.
    \item Déterminer une mesure $\theta$ de l'angle de $f$.
    \item Prouver que $f^{-1} = S_{(BO)} \circ S_{(BA)}$ où $S_{(BO)}$ et $S_{(BA)}$ désignent respectivement les symétries orthogonales d'axes $(BO)$ et $(BA)$, est la réciproque de $f$.
  \end{enumerate}
  \item Soit $g = S_{(OD)} \circ R_{(B,\frac{\pi}{2})}$, $R_{(B,\frac{\pi}{2})}$ désigne la rotation de centre $B$ et d'angle de mesure $\dfrac{\pi}{2}$.
  \begin{enumerate}[label=\alph*.]
    \item Justifier que $g$ est une symétrie glissée.
    \item Déterminer $g(B)$ et $g(C)$.
    \item En calculant $(g \circ g)(B)$, montrer que le vecteur de $g$ est $\vec{u} = \vv{BO}$.
    \item En déduire que l'axe de $g$ est la droite $(KO')$.
    \item Déterminer l'ensemble $(\Gamma)$ des points $M$ du plan tels que : $g(M) = f^{-1}(M)$.
  \end{enumerate}
  \item Soit $(H)$ l'hyperbole d'excentricité $e = 2$, de foyer $C$ et de directrice associée $(BA)$.
  \begin{enumerate}[label=\alph*.]
    \item Soit $J$ le point du plan, tel que $\vv{BJ} = \dfrac{1}{3}\vv{BC}$. Montrer que $G$ et $J$ sont les sommets de $(H)$.
    \item Placer le centre $\Omega$ de $(H)$.
    \item Tracer les asymptotes $(\Delta_1)$ et $(\Delta_2)$ de $(H)$.
    \item Tracer $(H_0)$, la branche de l'hyperbole dont le sommet est $J$.
    \item Le plan est muni d'un repère orthonormé direct $(B,\vv{BC},\vv{BA})$. En utilisant la définition monofocale de $(H)$ appliquée au foyer $C$, montrer qu'une équation cartésienne de $(H)$ dans ce repère est : $3x^2 + 2x - y^2 - 1 = 0$.
  \end{enumerate}
\end{enumerate}

% ====================== EXERCICE 3 ======================
\exo{3}{5 points}

\begin{enumerate}
  \item Résoudre l'équation différentielle $(E) : y'' + 3y' + 2y = 0$.
  \item Déterminer la solution particulière $h$ de $(E)$ qui admet en $x = 0$ un maximum égal à 1.
  \note{Modifier la question : « Déterminer la solution particulière $h$ de $(E)$ qui admet en $x = 0$ un extremum égal à 1. »}
  \item Soit $f$ la fonction numérique définie sur $\R$ par : $f(x) = -e^{-2x} + 2e^{-x}$. On désigne par $(C)$ sa courbe représentative dans le plan rapporté à un repère orthonormé $(O,\vi,\vj)$.
  \begin{enumerate}[label=\alph*.]
    \item Calculer $\displaystyle\lim_{x\to-\infty} f(x)$. On admet que $\displaystyle\lim_{x\to+\infty} f(x) = 0$.
    \item Déterminer la dérivée $f'$ de $f$.
    \item Étudier le sens de variation de $f$.
    \item Dresser le tableau de variation de $f$.
    \item Trouver le point d'intersection de la courbe $(C)$ avec l'axe des abscisses.
    \item Calculer $\displaystyle\lim_{x\to-\infty} \dfrac{f(x)}{x}$.
    \item Tracer $(C)$.
  \end{enumerate}
\end{enumerate}

% ====================== EXERCICE 4 ======================
\exo{4}{3 points}

On dispose de deux urnes $U_1$ et $U_2$. L'urne $U_1$ contient deux boules blanches et trois boules noires. L'urne $U_2$ contient deux boules blanches et deux boules noires. Dans chacune des urnes, les boules sont indiscernables au toucher.

On lance sur une table un dé cubique non truqué, dont les faces sont numérotées de 1 à 6.
\begin{itemize}
  \item Si le nombre apparu sur la face supérieure du dé est inférieur ou égal à 2, on tire une boule dans l'urne $U_1$ ;
  \item Si le nombre apparu sur la face supérieure du dé est strictement supérieur à 2, on tire une boule dans l'urne $U_2$.
\end{itemize}
On note $\mathcal{U}_1$ et $\mathcal{U}_2$ les événements suivants :
\begin{itemize}
  \item $\mathcal{U}_1$ : « le tirage s'effectue dans l'urne $U_1$ »
  \item $\mathcal{U}_2$ : « le tirage s'effectue dans l'urne $U_2$ »
  \item $\mathcal{B}$ : « Obtenir une boule blanche »
  \item $\mathcal{N}$ : « Obtenir une boule noire »
\end{itemize}

\begin{enumerate}
  \item \begin{enumerate}[label=\alph*.]
    \item Montrer que la probabilité de l'événement $\mathcal{U}_1$ est $\dfrac{1}{3}$.
    \item En déduire la probabilité de l'événement $\mathcal{U}_2$.
  \end{enumerate}
  \item Traduire par un arbre de probabilités, les données de l'énoncé.
  \item Montrer que la probabilité de tirer une boule noire est égale à $\dfrac{8}{15}$.
  \item On a tiré une boule noire. Calculer la probabilité qu'elle provienne de l'urne $\mathcal{U}_1$.
\end{enumerate}

\end{document}
